\begin{ouverture}[espace]{16}{Espace : sphère, boule et sections}{Partie 4 · Espace et géométrie}
{Solides · Volumes · Sphère et boule · Grands cercles · Sections planes}
\objectif{Reconnaître les solides et calculer leurs volumes.}
\objectif{Définir la sphère, la boule et leurs grands cercles.}
\objectif{Visualiser et calculer des sections planes.}
\prerequis{Aires usuelles, Pythagore, Thalès et unités de volume.}
\end{ouverture}
\pagestyle{otMain}\essentiel\fontsize{11}{15}\selectfont
\essSection{1}{Solides, sphère et boule}
\begin{propriete}Pour un pavé : $V=L\ell h$. Pour un prisme droit ou un cylindre : $V=\mathcal A_{\rm base}\times h$. Pour une pyramide ou un cône : $V=\mathcal A_{\rm base}\times h/3$. La hauteur est perpendiculaire à la base.\end{propriete}
\begin{definition}La sphère de centre O et de rayon R est l’ensemble des points de l’espace situés à la distance R de O. La boule de centre O et de rayon R est l’ensemble des points situés à une distance de O inférieure ou égale à R.\end{definition}
\begin{definition}Un grand cercle de la sphère est un cercle de centre O et de rayon R tracé sur la sphère. Deux points A et B de la sphère sont diamétralement opposés lorsque O est le milieu de [AB] ; alors $AB=2R$ et [AB] est un diamètre.\end{definition}
\begin{propriete}Le volume d’une boule de rayon $r$ est $V=\dfrac43\pi r^3$. Cette formule est admise.\end{propriete}
\begin{savoirfaire}[Assembler]Décomposer en solides usuels, additionner des volumes sans recouvrement ou soustraire un creux. Conserver $\pi$ jusqu’au résultat final, puis convertir : $1\text{ dm}^3=1\text{ L}$.\end{savoirfaire}
\clearpage
\essSection{2}{Sections planes}
\begin{propriete}La section d’un pavé par un plan parallèle à une face est un rectangle de mêmes dimensions que cette face ; par un plan parallèle à une arête, c’est un rectangle.\end{propriete}
\begin{propriete}La section d’un cylindre par un plan perpendiculaire à son axe est un disque de même rayon que la base ; par un plan parallèle à son axe, c’est un rectangle.\end{propriete}
\begin{propriete}La section d’une pyramide ou d’un cône par un plan parallèle à la base est une réduction de la base : même nature, dimensions proportionnelles. Le coefficient est le quotient de la distance au sommet par la hauteur.\end{propriete}
\begin{propriete}La section d’une boule par un plan est un disque. Si le plan est à la distance $d$ du centre, avec $d<R$, le rayon $r$ du disque vérifie $r^2=R^2-d^2$ d’après Pythagore. Le plan passant par le centre donne un grand cercle.\end{propriete}
\begin{savoirfaire}[Dessiner en vraie grandeur]Reporter les longueurs utiles à partir de la face avant de la perspective, ou les calculer, puis tracer la section à part, sans déformation.\end{savoirfaire}
\clearpage\rubrique{Découverte}
\begin{blocLivre}{ex-16-1}\begin{exercice}[1]{Un pavé}
Un pavé a pour dimensions 8 cm, 5 cm et 3 cm. Calculer son volume en cm$^3$ puis en litres.
\label{fig-ex-16-1}\figureGeo{espace-01}
\end{exercice}\end{blocLivre}
\begin{blocLivre}{ex-16-2}\begin{exercice}[1]{Un cylindre}
Un cylindre a un rayon de 3 cm et une hauteur de 10 cm. Donner son volume exact et au cm$^3$.
\label{fig-ex-16-2}\figureGeo{espace-02}
\end{exercice}\end{blocLivre}
\begin{blocLivre}{ex-16-3}\begin{exercice}[1]{Une pyramide}
Une pyramide a une base carrée de côté 6 cm et une hauteur de 9 cm. Calculer son volume. Quelle est la nature de sa base ?
\label{fig-ex-16-3}\figureGeo{espace-03}
\end{exercice}\end{blocLivre}
\begin{blocLivre}{ex-16-4}\begin{exercice}[1]{Un cône}
Un cône a pour diamètre 8 cm et hauteur 9 cm. Calculer son volume exact.
\label{fig-ex-16-4}\figureGeo{espace-04}
\end{exercice}\end{blocLivre}
\begin{blocLivre}{ex-16-5}\begin{exercice}[1]{Une boule}
Une boule a un rayon de 3 cm. Donner son volume exact puis arrondi au cm$^3$. Quelle est la longueur d’un grand cercle de sa sphère ?
\label{fig-ex-16-5}\figureGeo{espace-05}
\end{exercice}\end{blocLivre}

\clearpage\rubrique{Entraînement}
\begin{blocLivre}{ex-16-6}\begin{exercice}[2]{Sphère, boule et grands cercles}
La sphère de centre O a pour rayon 4 cm. A et B sont diamétralement opposés, C est sur la sphère, OM = 2 cm et ON = 6 cm. Quels points appartiennent à la sphère ? À la boule ? Calculer AB et la longueur du grand cercle passant par A et B.
\label{fig-ex-16-6}\figureGeo{espace-sphere-vocab}
\end{exercice}\end{blocLivre}
\begin{blocLivre}{ex-16-7}\begin{exercice}[2]{Une pyramide dans un cube}
ABCDEFGH est un cube d’arête 6 cm ; E est le sommet situé au-dessus de A. On considère la pyramide EABCD de sommet E et de base ABCD. Donner la nature de ses cinq faces, puis son volume. Quelle fraction du cube représente-t-elle ?
\label{fig-ex-16-7}\figureGeo{espace-pyramide-cube}
\end{exercice}\end{blocLivre}
\begin{blocLivre}{ex-16-8}\begin{exercice}[2]{Une section de pavé}
Un pavé de base 8 cm sur 5 cm et de hauteur 6 cm est coupé par un plan parallèle à la base à 2 cm de celle-ci. Quelle forme et quelle aire a la section ?
\label{fig-ex-16-8}\figureGeo{espace-section-pave}
\end{exercice}\end{blocLivre}
\begin{blocLivre}{ex-16-9}\begin{exercice}[2]{Sections d’un cylindre}
Un cylindre a un rayon de 3 cm et une hauteur de 10 cm. Décrire et calculer l’aire de sa section par un plan perpendiculaire à son axe, puis par un plan contenant son axe.
\label{fig-ex-16-9}\figureGeo{espace-section-cylindre}
\end{exercice}\end{blocLivre}
\begin{blocLivre}{ex-16-10}\begin{exercice}[2]{Une section de sphère}
Une sphère de centre O a un rayon de 5 cm. Un plan la coupe à une distance OH = 3 cm du centre ; H est le centre du cercle de section. Calculer le rayon de ce cercle et l’aire du disque de section.
\label{fig-ex-16-10}\figureGeo{espace-section-sphere}
\end{exercice}\end{blocLivre}

\clearpage\rubrique{Problèmes}
\begin{blocLivre}{ex-16-11}\begin{exercice}[3]{Un réservoir assemblé}
Un réservoir fermé est formé d’un cylindre de rayon 1 m et de hauteur 2 m, surmonté d’une demi-boule de même rayon. Calculer sa capacité au litre.
\label{fig-ex-16-11}\figureGeo{espace-reservoir}
\end{exercice}\end{blocLivre}
\begin{blocLivre}{ex-16-12}\begin{exercice}[3]{Une boule creusée}
Une boule pleine de rayon 5 cm contient une cavité sphérique de rayon 2 cm, entièrement intérieure. Quel volume de matière reste-t-il ?
\label{fig-ex-16-12}\figureGeo{espace-cavite}
\end{exercice}\end{blocLivre}
\begin{blocLivre}{ex-16-13}\begin{exercice}[3]{Un cône coupé}
Un cône de rayon 6 cm et de hauteur 12 cm est coupé par un plan parallèle à sa base, à mi-hauteur. Quelle est la section ? Calculer le volume du petit cône de même sommet, puis celui du tronc restant.
\label{fig-ex-16-13}\figureGeo{espace-cone-coupe}
\end{exercice}\end{blocLivre}
\begin{blocLivre}{ex-16-14}\begin{exercice}[3]{Une section de pyramide}
SABCD est une pyramide de base carrée ABCD de côté 6 cm et de hauteur SO = 9 cm. Un plan parallèle à la base coupe la pyramide à 3 cm du sommet S. Déterminer la nature et l’aire de la section, puis le volume de la petite pyramide obtenue.
\label{fig-ex-16-14}\figureGeo{espace-section-pyramide}
\end{exercice}\end{blocLivre}
\begin{blocLivre}{ex-16-15}\begin{exercice}[3]{Le cylindre d’Archimède}
Une boule de rayon 3 cm est posée dans un cylindre de rayon 3 cm et de hauteur 6 cm, qu’elle touche en haut, en bas et sur le côté. Calculer les deux volumes et le rapport du volume de la boule à celui du cylindre. Montrer que ce rapport ne dépend pas du rayon.
\label{fig-ex-16-15}\figureGeo{espace-archimede}
\end{exercice}\end{blocLivre}

\clearpage
